\section{Evaluation}
\label{sec:eval}

\subsection{Questions, Subject and Arms}

We ask five questions, all at fixture level.
\begin{description}
\item[RQ1] (a) Do the production audit and the ground truth agree under oracle
perception? (b) What does the exploration loop do: how many pairs does it execute,
what does replay cost, are revisit edges kept, and what happens when a
side-effecting step or a crash is injected?
\item[RQ2] How does the match between policy vocabulary and perception output
affect accuracy, and do the errors of a real channel and of perturbed perceivers
follow Proposition~\ref{prop:asym}?
\item[RQ3] Is the required-absent operation necessary? (Answered in the Table~\ref{tab:rq2} row ``$\mathcal R$ dropped'': an absence cannot intersect anything.)
\item[RQ4] Does $\phi$ keep volatile noise out, and which targeted edits does it
fail to separate?
\item[RQ5] On programmatic variants of the fixture, per mutant family, how do the
arms behave, and what do two hand-written invariant scripts, one reading the tag
encoding and one independent of it, achieve?
\end{description}

\paragraph{Subject and arms.}
\texttt{broken-admin} (Sect.~\ref{sec:example}) is a static six-page application
audited as role \texttt{guest}; its ground-truth files list four planted
violations (V-01 to V-04) and four named cycles. The policy has three forbidden
clauses and one required clause. V-03 lies three actions from the seed (index,
login, dashboard, admin). A violation is the triple (page, clause, clause type),
scored against the planted triples. Table~\ref{tab:arms} lists the arms. The
oracle arms use a fake perceiver that reads the \texttt{data-tag} annotations in
which the ground truth is itself written, so their agreement concerns clause logic
and encoding, not perception. No arm uses a browser, Neo4j or Redis, and the AX
tree is used only in the fingerprint tests.

\begin{table}[tbp]
\caption{Experimental arms. \emph{Prod.} is the production negation engine;
\emph{fake} is the in-memory fake's own set algebra. RQ4 uses E2 and E8.}
\label{tab:arms}
\setlength{\tabcolsep}{2.5pt}
\footnotesize
\begin{tabular}{@{}llll@{}}
\toprule
Arm & Perception & Audit & RQ: JSON key \\
\midrule
A1 oracle, fake algebra & fake, reads \texttt{data-tag} & fake & 1a, 3: E1\_E3 \\
A2 oracle, production & fake, reads \texttt{data-tag} & prod. & 1a,3: E6; 5: E9 \\
A3 DOM channel & rule-based builder & prod. & 2: E1\_E3; 5: E9 \\
A4 DOM + lexical & A3, hyphen as space & patched & 2: E1\_E3 \\
A5 perturbed oracle & A2 maps, perturbed & prod. & 2: E7 \\
A6 tag script & HTML parse of \texttt{data-tag} & hand-written & 5: E9 \\
A7 loop & fake, scripted crawler & fake & 1b: E5, E11 \\
A8 XPath script & element types, text & hand-written & 5: E10, E12--E14 \\
\bottomrule
\end{tabular}
\end{table}

\paragraph{Environment and artifact.}
The script \path{experiments/run_experiments.py} writes
\path{data/experiments.json} (both under \path{docs/report/}), which records Python
3.13.13, LangGraph 1.2.10 and repository commit \texttt{4487574} (the code
under \texttt{apps/}, \texttt{tests/} and \texttt{test-apps/} was unmodified, but the
stored flag \texttt{code\_tree\_clean\_vs\_head} is false because the E14 fixtures,
an untracked folder under \texttt{tests/}, exist) and a map from each RQ to its keys. Seed 7 drives all
randomised parts. The script, data file and E14 fixtures are not committed and no archived
release exists.

\subsection{RQ1: Audit Agreement and Loop Behaviour}
\label{sec:rq1}

\paragraph{(a) Agreement under oracle perception.}
Both oracle arms report exactly the four planted triples (P, R and F1 of 1.00;
Table~\ref{tab:rq2}, rows 1--2) and agree page by page. Arm A2 sends the fake's
semantic maps through the production \texttt{negation.audit.audit}, so the
production operators are exercised. It remains a consistency check: both arms read the ground truth's tag encoding, so
a perfect score is close to guaranteed, and it is neither a perception measurement
nor an end-to-end result.

\paragraph{(b) The exploration loop.}
We ran the LangGraph loop over a scripted crawler built from \texttt{app.json}
with the fake perceiver and graph (A7); the framework-free reference loop gave
identical counters and graph. Table~\ref{tab:loop} gives the counters.
\begin{table}[tbp]
\caption{Loop counters (A7); the 25 shuffles give identical counts. ``Dup.'': pairs
executed as a frontier action more than once; ``replay/steps'': re-navigations and
re-fired actions; ``revisits'': revisit edges. In the last row the
login-to-dashboard step is a side-effecting action.}
\label{tab:loop}
\setlength{\tabcolsep}{3.5pt}
\footnotesize
\begin{tabular}{@{}lcccccc@{}}
\toprule
Scenario & States & Pairs & Dup. & Replay/steps & Revisits & Skipped \\
\midrule
Uninterrupted; 25 link orders & 6 & 11 & 0 & 9/21 & 6 & 0 \\
Login step \texttt{submit} or \texttt{fill} & 4 & 4 & 0 & 2/3 & 1 & 3 \\
\bottomrule
\end{tabular}
\end{table}

The uninterrupted run ends by frontier exhaustion, reaches all 6 states, persists
all 11 transitions (5 tree edges, 6 revisit edges), finds all four planted
violations and executes no pair twice as a frontier action; violation sets are
also identical across link orders. Replay is nonetheless the dominant cost: 21
re-fired actions for 11 frontier executions (1.91 per pair; the bound of
Proposition~\ref{prop:term} is $11\cdot3=33$, with $D=3$ the observed maximum
depth and a configured limit of 10), and \emph{index $\to$ login} is executed 10
times, once as a frontier action and nine times as a replay prefix. With
\texttt{click} counted as replay-safe, that is harmless only if clicks are
navigational.

Making the login step a \texttt{submit} or \texttt{fill} shows the other side of
the rule: every candidate behind it has an unreplayable path and is skipped with a
recorded reason, so 4 states are reached, 3 actions are skipped and one of four
planted violations (V-01) is found. The rule trades coverage for safety.

\paragraph{Crash and resume.}
Proposition~\ref{prop:term} excludes crashes. We injected an exception in the
\texttt{act} node right after a claim, and separately in the edge write, at each of
the 11 claims (13 on the variant below), using LangGraph's in-memory checkpointer, and resumed in three ways:
from the checkpoint with the same fake browser, from it with a fresh crawler that
has opened no page, and with no checkpoint, only the persisted claim set, as a new
run on the same graph (Table~\ref{tab:resume}). The fake action fails before it has
any effect, so we count \emph{attempts}; whether an attempt had an effect in a real
AUT is unknowable here.
\begin{table}[tbp]
\caption{Simulated crash and resume at 11 positions. ``Twice'': positions at which
a claimed pair was attempted again after resume. ``Graph equal'': positions at which
the final graph equals the uninterrupted run's (a shortfall means exploration was
lost).}
\label{tab:resume}
\setlength{\tabcolsep}{3.5pt}
\footnotesize
\begin{tabular}{@{}llcc@{}}
\toprule
Resume mechanism & Crash point & Twice & Graph equal \\
\midrule
Checkpoint, same browser & action, after claim & 11/11 & 11/11 \\
Checkpoint, fresh browser & action, after claim & 11/11 & 9/11 \\
Checkpoint, same browser & edge write & 0/11 & 11/11 \\
Checkpoint, fresh browser & edge write & 0/11 & 10/11 \\
Claim set only, new run & action, after claim & 0/11 & 0/11 \\
\bottomrule
\end{tabular}
\end{table}

LangGraph documents that a resumed node restarts from its beginning and that side
effects should be idempotent or wrapped in tasks \cite{langgraph2026durable}, so the
mechanism is the framework's documented behaviour. The finding is the gap between
it and the claim handling intended in Sect.~\ref{sec:model}: resuming the
checkpoint re-runs the interrupted \texttt{act} node, which does not consult the
claim set again, so the in-flight pair is attempted twice at all 11 positions; a
crash after the action but before the edge is written is handled correctly. A
browser that is not where the checkpoint says it is leaves the graph incomplete in
3 of 22 cases. Restarting from the claim set alone never repeats a claimed pair
but never continues either: every action of the seed is already claimed, so the
new run claims no pair, visits one state, and leaves 0 to 10 of the 11 transitions
persisted. That last result is specific to this fixture, whose seed has one
action. On a variant with a second seed action (an \texttt{About} page; 7 states, 13
pairs; E11) the checkpoint results are the same (13/13 attempted twice after a
claim, 0/13 after an edge write), but a claim-set restart now continues through the
seed action that was not yet claimed: with a crash at the first claim the new run
executes 2 pairs, and from the second claim on it executes none, leaving 1 to 12
of 13 transitions persisted. The positions share one mechanism and are not samples. Claim-before-execute thus gives at-most-once only in the
second mode, and only by giving up completeness.

\subsection{RQ2: Vocabulary Match, Perception Error and Error Direction}
\label{sec:rq2}

\begin{table}[tbp]
\caption{Accuracy on \texttt{broken-admin} by arm, with errors split by clause type
($\mathcal F$: forbidden-present, $\mathcal R$: required-absent). Genuine TP excludes
a required-absent hit when the clause also fires on pages that carry the element:
it then fires everywhere and the hit is a coincidence.}
\label{tab:rq2}
\setlength{\tabcolsep}{2.2pt}
\footnotesize
\begin{tabular}{@{}lccccccccc@{}}
\toprule
Arm & P & R & F1 & Gen.\ TP & FN$_\mathcal F$ & FP$_\mathcal F$ & FN$_\mathcal R$ & FP$_\mathcal R$\\
\midrule
A1 oracle, fake algebra & 1.00 & 1.00 & 1.00 & 4 & 0 & 0 & 0 & 0 \\
A2 oracle, production & 1.00 & 1.00 & 1.00 & 4 & 0 & 0 & 0 & 0 \\
A1, $\mathcal R$ dropped (RQ3: loses V-04) & 1.00 & 0.75 & 0.86 & 3 & 0 & 0 & 1 & 0 \\
A3 DOM channel & 0.29 & 0.50 & 0.36 & 1 & 2 & 0 & 0 & 5 \\
A4 DOM + lexical & 1.00 & 0.50 & 0.67 & 2 & 2 & 0 & 0 & 0 \\
\bottomrule
\end{tabular}
\end{table}

\paragraph{A vocabulary-mismatch case study.}
Arm A3 builds the semantic map from the HTML with rule-based tagging
(Table~\ref{tab:rq2}). Its precision of 0.29 and recall of 0.50 do not measure
perception quality; they follow from an inconsistency inside the prototype. The
policy subjects are \texttt{admin-link}, \texttt{debug-panel} and
\texttt{skip-to-content}, but the tagger emits shorter tags for the admin link and
the debug section and none equal to the skip subject. V-01 and V-02 are missed,
and the skip clause fires on all six pages: five false positives and one
\emph{coincidental} hit on V-04. Only V-03 is detected for the right reason, via
the delete-user capability, so genuine recall is 0.25, not 0.50.

Arm A4 normalises hyphens inside the matcher (in-process patch). It was designed
after seeing A3's errors, on the same four items, so it is an exploratory
diagnosis, not a test. It removes the five false positives, so V-04 is now detected
for the right reason; V-01 and V-02 stay missed, because deciding that a link
labelled ``Admin'' is an \texttt{admin-link} is a semantic judgement. Whether to rename the policy subjects or let perception read the annotations is an
unresolved design decision, on which the end-to-end result depends.

\paragraph{A falsifiable test of the error direction.}
A3 and A4 err in the direction that Proposition~\ref{prop:asym} reads off for
under-reporting, but since a vocabulary mismatch counts as under-reporting by
definition, they cannot refute it. We therefore perturbed the oracle maps (A5):
under-reporting drops each policy-relevant element with probability $p$,
over-reporting adds each absent policy subject as a spurious element with
probability $p$, and the production audit runs on every perturbed state
(200 trials for $p=0.25$ and $0.5$, one deterministic run for $p=1$; E7). For example, at $p=0.5$ under-reporting gives on average 1.57 forbidden misses and 2.58 required false alarms, over-reporting 7.79 forbidden false alarms and 0.47 required misses.
In all 802 runs, under-reporting produced no forbidden false positive and no
required false negative, and over-reporting produced no forbidden false negative
and no required false positive. This confirms that the engine is monotone and antitone as Eqs.~(\ref{eq:forb}) and~(\ref{eq:req}) require. It says nothing about how a real VLM errs.

\subsection{RQ4: Fingerprint Noise and Targeted Edits}
\label{sec:fp}

Each page is re-rendered 50 times (seed 7) with fresh session and CSRF attributes,
a generated id, a request id, a 13-digit timestamp and a CDP-shaped AX tree with new
node ids. A naive hash of URL, DOM and AX JSON gives 300 distinct ids for six
logical pages, $\phi$ gives 6. This is a regression check against noise $N$ was written to remove.

The informative part is twelve targeted edits to the \texttt{admin} page. Eight
change the id (removing or adding a button, renaming a label, swapping links,
changing the URL query or an AX name, and, debatably, adding a \texttt{hidden}
element or editing a comment). Four do not: a whitespace change (by design) and three content
changes that fall inside the normaliser's patterns, namely a visible 10-digit
number, a visible word with a \texttt{req\_} prefix, and an identifier-like
\texttt{id} value (\texttt{user-1234}). Those three are over-merges: states that
differ only in such content are audited once, which matters when a policy subject
depends on that content (not the case here), and the patterns are broad. The
hidden-element and comment edits show the opposite risk of over-splitting on pages
with hidden or generated markup. We did not measure how often either occurs on
real pages.


\subsection{RQ5: Mutants, Baselines and Further Apps}
\label{sec:rq5}

We generated variants of \texttt{broken-admin} in memory (nothing is written to
the test applications): all 24 applicable single operations (remove or add
forbidden elements, remove or restore the skip link), 40 seeded compound mutants,
and label renames that keep the defects but change visible text. A generator
ledger gives each variant's ground truth independently of the audit code. Two
hand-written scripts, our own re-implementations and neither Crawljax nor ATUSA,
serve as baselines. A6 reads \texttt{data-tag}, like the oracle arms. A8 does not:
it has four invariants in the XPath subset of Python's ElementTree over element
types, link targets and visible text (for instance, an \texttt{a} element whose
text is ``Admin''), plus a 26-line shared HTML-to-tree helper. The authors knew the
labels when writing it; it was written before its first run on stripped pages,
which we state but cannot evidence, as the artifact has no history. Every variant
is also run on a copy with the \texttt{data-tag} attributes removed.

\begin{table}[tbp]
\caption{Errors per mutant family, as FN$_\mathcal F$/FP$_\mathcal R$ summed over the
family's $n$ variants (no arm produced FP$_\mathcal F$ or FN$_\mathcal R$ anywhere).
``Planted'': forbidden/required violations in the family. A2 and A6 on intact
tags make no error in any family, and A8 on stripped copies equals A8 on intact
ones. A2s is the oracle perceiver run without its annotations (stripped copies),
shown as a reference and not as a competitor. The label-rename row includes the
combined variant, so its A8 totals double-count the singles.}
\label{tab:mutants}
\setlength{\tabcolsep}{3pt}
\footnotesize
\begin{tabular}{@{}lcccccc@{}}
\toprule
Family & $n$ & Planted & A3 & A2s & A8 \\
\midrule
Remove planted forbidden & 3 & 6/3 & 4/15 & 6/15 & 0/0 \\
Add forbidden element & 15 & 60/15 & 40/75 & 60/75 & 0/0 \\
Remove or restore skip link & 6 & 18/10 & 12/26 & 18/26 & 0/0 \\
Compound & 40 & 162/47 & 113/193 & 162/193 & 0/0 \\
Label rename & 4 & 12/4 & 8/20 & 12/20 & 4/10 \\
\bottomrule
\end{tabular}
\end{table}

On the 64 shape-preserving variants A2, A6 and A8 report exactly the ledger and A3
matches it on none. A2 and A6 agree because both read the tag spelling in which the
mutants are written, so 64/64 shows consistency, not advantage. On unannotated pages
the informative comparison is A3 against A8: a short label-based script beats the
shipped DOM channel (stripped original fixture: A3 2 true and 5 false positives,
A8 4 and 0), because A3's vocabulary does not match the policy subjects
(Sect.~\ref{sec:rq2}). A8's exactness is not independent evidence, since the
mutator inserts the label texts A8 was written against. When visible text changes
it loses: of three independent renames, ``Admin'' and ``Delete user'' each cost A8
one missed forbidden violation and ``Skip navigation'' makes the skip clause fire
on five pages, while A2 and A6 are unaffected; tag-spelling renames (E9) are the
mirror image. Each encoding fails where the other does not. The 68 variants
rearrange the same six pages, so their counts are not independent samples.

\paragraph{Three further apps.}
The repository holds three more apps from the same project,
\texttt{broken-rbac-demo}, \texttt{missing-consent-demo} and
\texttt{broken-feature-flag-demo}: one to four static, unannotated pages and one
documented defect each; we read them and modified nothing. Clauses and
ground-truth triples were derived from each README sentence first (the ordering is
an authorial statement the artifact cannot evidence), and the XPath invariants were
written after reading the scored HTML, so A8's results are in-sample and say nothing
about generalisation (E12). A8 finds all three defects with no false positive. A3,
with README-derived subjects, finds only the consent one, vacuously: the clause
fires on any page, since the tagger has no consent vocabulary, and still fires when
a banner is inserted in memory; the consent row therefore does not discriminate for
A8 either. A3 misses the other two for the vocabulary reason of RQ2; with subjects
respelled in the tagger's vocabulary (post hoc) it finds both but also flags
\texttt{admin.html}. The truth is narrower than the README text: the feature-flag
policy sentence also names an analytics panel, which A8's exact-heading query
misses (under a truth that includes it, A8 finds 1 of 2 and A3 0 of 2), and the RBAC
dashboard leaks a delete button no clause covered.

\paragraph{Four generated apps (E14).}
Four apps (\texttt{shop-hub}, \texttt{med-clinic}, \texttt{hr-suite},
\texttt{config-portal}; 7 unannotated pages each, three violations and clean-page
lists each) were generated by an AI tool within the same project. They add variety
and a more blind protocol, not independence. The A8 invariants were written from the
policy files only, before opening any page, ground truth or README, and saved with
their SHA-256 and a UTC timestamp (E14). The order is self-attested: the timestamp is
a file time, all fixtures existed before it, and the hash shows only that the file
was not edited afterwards. A3 uses the policy subjects verbatim, and nothing was
changed after the results.

The fixtures are weak. The ground truth is the generator's own assertion and is
unvalidated. Six of the eight forbidden-violation pages have no inbound link, login
pages link nowhere, and roles are nominal in static pages. Several clean-page labels
contradict the policy wording or the HTML: the shop-hub admin clause speaks of a
\emph{link} that exists nowhere; shop-3 holds app-wide through the Support link yet
flags only the home page; med-3 lists \texttt{appointments.html} as clean although
it lacks emergency information; cfg-3 lists \texttt{settings.html} as clean although
it has an organisation field and a save button; and hr-3 flags the benefits page
although the dashboard shows the status. False-positive counts are therefore only as
meaningful as these lists.

A3 finds none of the 8 forbidden violations. On required clauses it fires on every
page, so its 4 hits are coincidental, with 19 false positives on listed clean pages
and 33 hits on unlisted pages (neither truth nor clean, so not counted as false
positives), because its vocabulary matches none of the subjects. A8 finds 3 of 8
forbidden violations (checkout, staff-portal and team pages) with no false positive;
it misses five, where the violation is a whole page reached without a link in the
form it guessed or the label differs from its guess (``Run Payroll Calculation'').
Its raw required tally is 3 of 4 with 4 false positives and 12 unlisted hits, but
applying Table~\ref{tab:rq2}'s genuine-TP rule as for A3 removes two of the three
hits (the clause also fires on listed pages that carry the element), and the third,
on \texttt{config-portal}, is vacuous: no fixture page has an \texttt{id} attribute,
so its \texttt{id} queries can never match. A8's genuine required recall is thus 0 of
4 to 1 of 4, depending on how ``carry'' is read, and it misses the benefits-status
violation because it tests for a link, not the status display. A8's invariants also
carry page scope (``the home page'', ``not the login page'') that the engine lacks,
which A3 cannot express. All of A8's false positives and 7 of A3's fall on the three
double-negation clauses (``must not fail to find'', ``should never be unable to
find''), encoded as required presence as the policy file does. A hand-written label
script thus helps only where its author guessed the labels, a more modest result
than on \texttt{broken-admin}.

\subsection{Threats to Validity}
\label{sec:threats}

This is the single statement of the evidence's limits. \emph{External.} The
evidence is one small, static, hand-built application with four planted defects
and its variants, plus three tiny single-defect apps and four generated apps (untracked fixtures of weak soundness), all from the same project; it
shows that the mechanisms behave as designed on these and is not independent
validation. No genuinely independent subject and no real Crawljax-plus-ATUSA
baseline were available; A8 is a partial, self-written substitute.
\emph{Construct.} The oracle arms read the ground truth's own encoding; A3 and A4
work on HTML with no AX tree; the loop experiments use fakes, and the crash test
injects in-process exceptions with an in-memory checkpointer and a fake browser,
so it probes loop logic and not Redis, Neo4j or Playwright behaviour. No
browser-driven, VLM or hybrid result is reported. \emph{Internal.} The authors
wrote the application, policy, tool, mutation operators and both baselines, knew
the labels A8 queries on \texttt{broken-admin} (not on the E14 apps), and chose the rename variants; A4 was designed after
inspecting A3's errors. \emph{Statistical.} Runs are deterministic or use one seed, without
confidence intervals; positive counts are small.
